% === Section 9: Conclusions ===
\section{Conclusions}
\label{sec:conclusions}\label{sec:conclusion}

We presented a PRISMA-compliant systematic review of foundation models for
retail demand forecasting, augmented with original gap-filling experiments that
fill a structural hole in the literature: no existing paper reports both FM and
gradient-boosted tree results on the same retail dataset under matched
conditions (\S\ref{sec:literature}).  The study synthesised 185 extraction rows
from 38 studies (2020--2026) with 45 experiment rows across M5, Favorita, and
Rohlik~v2 under a unified rolling-origin protocol.  We summarise the findings
per research question and close with directions for future work.

\subsection{Answers to the research questions}
\label{subsec:answers-rq}

\paragraph{RQ1 (Accuracy): Do zero-shot foundation models outperform
gradient-boosted tree baselines on retail \WAPE{}?}

No, not in general.  The cross-paper pooled meta-regression on $k = 10$
bucket-level deltas finds $\Dhat(\text{FM} - \text{ML\_TREE}) = +0.101$
\WAPE{} ($p = 0.505$) on the full pool --- a sign flip driven by proper
precision-weighting of the M5 \WRMSSE{} bucket --- and $\Dhat = -0.044$
\WAPE{} ($p = 0.17$, $Q$ $p = 0.85$) after excluding M5.  The confidence
interval on the excl-M5 estimate $[-0.11, 0.03]$ includes zero with negligible
heterogeneity.  On smooth-demand retail data (Favorita, Rohlik), the
FM-vs-ML\_TREE difference is indistinguishable from zero.

The M5 \WAPE{} cells ($\Delta \approx -0.54$ to $-0.78$) show a large
apparent FM advantage, but this reverses on \WRMSSE{} ($\Delta = +0.41$)
because per-series \WAPE{} on intermittent demand penalises tree models via
zero-denominator explosions.  M5 is a dataset-specific finding driven by metric
sensitivity, not a family-level effect.

\paragraph{RQ2 (Moderators): Which data characteristics moderate the
FM-vs-ML\_TREE gap?}

Demand intermittency (H1) is the primary moderator: M5 (zero-day fraction
${\sim}$70\%) is the only dataset where the FM-vs-ML\_TREE sign depends on
metric choice, and it is the only dataset that generates significant
heterogeneity in the pool.  Covariate richness (H2) is direction-consistent:
on Favorita, where covariates carry real signal, \texttt{lightgbm\_cov} matches
the FM within 0.3~pp.  The LightGBM multi-step protocol (H3) is conditional on
intermittency: direct wins on M5 (intermittent) by 17--22\%, recursive wins on
Favorita and Rohlik (smooth) by 3--10\%, vindicating the \S\ref{sec:experiments}
conditional synthesis.  The interaction term H4 is untestable at $k = 10$.

All moderator CIs are wide at the current pool size; we report
direction-consistency rather than CI-based significance tests and flag this as
the primary limitation of the current analysis.

\paragraph{RQ3 (Cost): What is the cost-accuracy Pareto frontier?}

Consumer-hardware FM inference (MacBook M-series MPS) costs ${\sim}$\$0.003
per 1,000 series per horizon --- $16\times$ cheaper than Azure E4DS\_V4 batch
baselines in USD.  However, the Polish-grid CO$_2$ footprint is $44\times$
higher per \$ than the Swedish-hydro Azure grid, and wall time is $4\times$
slower.  On smooth-demand data where FM and \LGBM{} accuracy are within 1~pp
\WAPE{}, the cost differential favours the FM when no cloud budget exists and
favours \LGBM{} when cloud infrastructure is already provisioned.

\paragraph{RQ4 (Guidance): When should a retail team deploy a foundation
model?}

The \S\ref{sec:framework} decision framework distils the findings into three
conditions where the FM pays off:

\begin{enumerate}
  \item \textbf{No feature engineering capacity.}  When the team cannot build
    or maintain a covariate pipeline, zero-shot FM inference on consumer
    hardware delivers competitive accuracy ($<$1~pp \WAPE{} gap on
    smooth-demand data) with zero data preparation.

  \item \textbf{Cold-start series.}  For new SKUs with no sales history,
    zero-shot FMs produce reasonable forecasts from day one, while global
    LightGBM requires at least a partial training window.

  \item \textbf{Aggregate-level forecasting on intermittent data.}  When the
    business KPI is at the category or store level (not per-SKU), the aggregate
    \WAPE{} of FMs on M5 (${\sim}$0.85) is comparable to \LGBM{}, and the FM
    avoids the per-series \WAPE{} pathology entirely.
\end{enumerate}

The FM does NOT pay off when: (a)~the team has rich covariates and the capacity
to use them (\LGBM{} matches or beats FM), (b)~per-series accuracy on
intermittent demand matters (\WRMSSE{} favours \LGBM{} by 45~pp on M5), or
(c)~batch throughput at scale is the constraint (\LGBM{} on CPU is faster per
series than FM on MPS for $>$10K series).

\subsection{Future work}
\label{subsec:future-work}

\paragraph{Covariate-aware FM evaluation.}
Chronos-2 (v2) and Moirai~2.0 support exogenous covariates.  Running these
models with promotions and calendar features on Favorita and Rohlik would test
whether the covariate channel closes the FM-vs-\LGBM{} gap on smooth-demand
data.  This is the single highest-value extension of the current work.

\paragraph{Larger cross-paper pool.}
The $k = 10$ pool limits moderator analysis to direction-consistency.  As more
FM papers report per-series \WAPE{} on individual retail datasets under
rolling-origin protocols, the pool will grow and enable CI-based moderator
tests.  We encourage the community to report per-series and aggregate \WAPE{}
alongside skill scores, and to include at least one tree-based baseline in
every retail FM evaluation.

\paragraph{Probabilistic metrics.}
CRPS, quantile loss, and calibration error would complement the point-forecast
\WAPE{} analysis and may reveal FM advantages in the tails of the predictive
distribution that \WAPE{} cannot capture.

\paragraph{Fine-tuned FM evaluation.}
Our study is restricted to zero-shot FMs.  Lightweight fine-tuning (e.g., LoRA
on Chronos-2 with 1\% of target data) could shift the FM-vs-ML\_TREE balance,
especially on covariate-rich datasets where the FM's pre-trained features are
augmented with task-specific signal.

\paragraph{Beyond retail.}
The conditional finding (intermittency $\times$ metric choice drives the
FM-vs-ML\_TREE comparison) should be tested on spare-parts, pharmaceutical, and
e-commerce demand datasets, which span the full intermittency spectrum.
